% manuscript-sha256: edeb29b681854d6ff67bdf072e4d3163995e068973662ebb77220ba009345238
% Options for packages loaded elsewhere
\PassOptionsToPackage{unicode}{hyperref}
\PassOptionsToPackage{hyphens}{url}
\documentclass[
]{article}
\usepackage{xcolor}
\usepackage[margin=1in]{geometry}
\usepackage{amsmath,amssymb}
\setcounter{secnumdepth}{-\maxdimen} % remove section numbering
\usepackage{iftex}
\ifPDFTeX
  \usepackage[T1]{fontenc}
  \usepackage[utf8]{inputenc}
  \usepackage{textcomp} % provide euro and other symbols
\else % if luatex or xetex
  \usepackage{unicode-math} % this also loads fontspec
  \defaultfontfeatures{Scale=MatchLowercase}
  \defaultfontfeatures[\rmfamily]{Ligatures=TeX,Scale=1}
\fi
\usepackage{lmodern}
\ifPDFTeX\else
  % xetex/luatex font selection
  \setmainfont[]{Times New Roman}
  \setmonofont[]{Consolas}
  \setmathfont[]{Latin Modern Math}
\fi
% Use upquote if available, for straight quotes in verbatim environments
\IfFileExists{upquote.sty}{\usepackage{upquote}}{}
\IfFileExists{microtype.sty}{% use microtype if available
  \usepackage[]{microtype}
  \UseMicrotypeSet[protrusion]{basicmath} % disable protrusion for tt fonts
}{}
\makeatletter
\@ifundefined{KOMAClassName}{% if non-KOMA class
  \IfFileExists{parskip.sty}{%
    \usepackage{parskip}
  }{% else
    \setlength{\parindent}{0pt}
    \setlength{\parskip}{6pt plus 2pt minus 1pt}}
}{% if KOMA class
  \KOMAoptions{parskip=half}}
\makeatother
\usepackage{color}
\usepackage{fancyvrb}
\newcommand{\VerbBar}{|}
\newcommand{\VERB}{\Verb[commandchars=\\\{\}]}
\DefineVerbatimEnvironment{Highlighting}{Verbatim}{commandchars=\\\{\}}
% Add ',fontsize=\small' for more characters per line
\newenvironment{Shaded}{}{}
\newcommand{\AlertTok}[1]{\textcolor[rgb]{1.00,0.00,0.00}{\textbf{#1}}}
\newcommand{\AnnotationTok}[1]{\textcolor[rgb]{0.38,0.63,0.69}{\textbf{\textit{#1}}}}
\newcommand{\AttributeTok}[1]{\textcolor[rgb]{0.49,0.56,0.16}{#1}}
\newcommand{\BaseNTok}[1]{\textcolor[rgb]{0.25,0.63,0.44}{#1}}
\newcommand{\BuiltInTok}[1]{\textcolor[rgb]{0.00,0.50,0.00}{#1}}
\newcommand{\CharTok}[1]{\textcolor[rgb]{0.25,0.44,0.63}{#1}}
\newcommand{\CommentTok}[1]{\textcolor[rgb]{0.38,0.63,0.69}{\textit{#1}}}
\newcommand{\CommentVarTok}[1]{\textcolor[rgb]{0.38,0.63,0.69}{\textbf{\textit{#1}}}}
\newcommand{\ConstantTok}[1]{\textcolor[rgb]{0.53,0.00,0.00}{#1}}
\newcommand{\ControlFlowTok}[1]{\textcolor[rgb]{0.00,0.44,0.13}{\textbf{#1}}}
\newcommand{\DataTypeTok}[1]{\textcolor[rgb]{0.56,0.13,0.00}{#1}}
\newcommand{\DecValTok}[1]{\textcolor[rgb]{0.25,0.63,0.44}{#1}}
\newcommand{\DocumentationTok}[1]{\textcolor[rgb]{0.73,0.13,0.13}{\textit{#1}}}
\newcommand{\ErrorTok}[1]{\textcolor[rgb]{1.00,0.00,0.00}{\textbf{#1}}}
\newcommand{\ExtensionTok}[1]{#1}
\newcommand{\FloatTok}[1]{\textcolor[rgb]{0.25,0.63,0.44}{#1}}
\newcommand{\FunctionTok}[1]{\textcolor[rgb]{0.02,0.16,0.49}{#1}}
\newcommand{\ImportTok}[1]{\textcolor[rgb]{0.00,0.50,0.00}{\textbf{#1}}}
\newcommand{\InformationTok}[1]{\textcolor[rgb]{0.38,0.63,0.69}{\textbf{\textit{#1}}}}
\newcommand{\KeywordTok}[1]{\textcolor[rgb]{0.00,0.44,0.13}{\textbf{#1}}}
\newcommand{\NormalTok}[1]{#1}
\newcommand{\OperatorTok}[1]{\textcolor[rgb]{0.40,0.40,0.40}{#1}}
\newcommand{\OtherTok}[1]{\textcolor[rgb]{0.00,0.44,0.13}{#1}}
\newcommand{\PreprocessorTok}[1]{\textcolor[rgb]{0.74,0.48,0.00}{#1}}
\newcommand{\RegionMarkerTok}[1]{#1}
\newcommand{\SpecialCharTok}[1]{\textcolor[rgb]{0.25,0.44,0.63}{#1}}
\newcommand{\SpecialStringTok}[1]{\textcolor[rgb]{0.73,0.40,0.53}{#1}}
\newcommand{\StringTok}[1]{\textcolor[rgb]{0.25,0.44,0.63}{#1}}
\newcommand{\VariableTok}[1]{\textcolor[rgb]{0.10,0.09,0.49}{#1}}
\newcommand{\VerbatimStringTok}[1]{\textcolor[rgb]{0.25,0.44,0.63}{#1}}
\newcommand{\WarningTok}[1]{\textcolor[rgb]{0.38,0.63,0.69}{\textbf{\textit{#1}}}}
\usepackage{longtable,booktabs,array}
\usepackage{calc} % for calculating minipage widths
% Correct order of tables after \paragraph or \subparagraph
\usepackage{etoolbox}
\makeatletter
\patchcmd\longtable{\par}{\if@noskipsec\mbox{}\fi\par}{}{}
\makeatother
% Allow footnotes in longtable head/foot
\IfFileExists{footnotehyper.sty}{\usepackage{footnotehyper}}{\usepackage{footnote}}
\makesavenoteenv{longtable}
\setlength{\emergencystretch}{3em} % prevent overfull lines
\providecommand{\tightlist}{%
  \setlength{\itemsep}{0pt}\setlength{\parskip}{0pt}}
\usepackage{bookmark}
\IfFileExists{xurl.sty}{\usepackage{xurl}}{} % add URL line breaks if available
\urlstyle{same}
\hypersetup{
  hidelinks,
  pdfcreator={LaTeX via pandoc}}

\author{}
\date{}

\allowdisplaybreaks
\begin{document}

\begin{center}
{\LARGE\bfseries Exact reductions and quantitative bounds for mutually unbiased bases in dimension six\par}
Jeremy H. Carroll\par
2026-10-16\par
Version 1.1.1\par
\end{center}

\subsection{Abstract}\label{abstract}

This technical checkpoint separates finite exact calculations, reported
Lean implications, and unresolved links to concrete mutually unbiased
bases. A pinned replay of the finite-corner classification, together
with concrete normalization and transport, now reduces actual quartet
existence to completion of a retained reconstruction. A separate exact
Cauchy--Binet calculation gives a minor with squared determinant at
least \(54/5\) and inverse-entry squared moduli at most \(10/27\) in
every raw Hadamard. Whether nonextension holds for every raw Hadamard
remains open. The bounded minor now also supplies a complete block
reconstruction chart. Clearing its nonzero determinant denominator
yields explicit division-free matrix equations equivalent to actual
quartet existence, with both additional bases retained. A checked affine
completion of the balanced transition matrices gives an equivalent
system of 648 rational polynomial equations of degree at most two in 320
bounded real coordinates. Its infeasibility is not established. The code
also contains cyclotomic fixtures, rational interval routines,
conditional atlas and gluing interfaces, and finite parity calculations.
The Route C obstruction still needs a certified Fourier-triple landing
premise. The declared 56-cell cover still needs analytic curvature
certificates. The operator interface admits a zero-action countermodel
and therefore needs stronger mathematical definitions before it can
express the intended exclusion theorem. Separate concrete constructions
now supply the continuous two-circulant family, variable-basis
equations, and an actual fourth-basis candidate from normalized positive
flat compatible moments. The needed empty varieties and full analytic
cover remain open. No unconditional dimension-wide MUB bound follows
from this package at this checkpoint.

\subsection{1. Mathematical problem and
normalization}\label{mathematical-problem-and-normalization}

This article presents reductions and quantitative bounds, rather than a
solution of the dimension-six MUB problem. Its contribution is a
separation of concrete matrix reductions, finite calculations,
conditional interfaces and remaining exclusion obligations. Section 2.1
compares the reductions with existing family results and the
finite-corner classification.

Formal-verification statements below describe the source reports
identified in Section 6. The standalone preprint archive contains the
manuscript and rendering machinery; it contains no Lean project or
computational research supplement. Those formal statements have not
received a complete fresh replay for this release candidate and are not
certified by its checksum verifier.

Let \(B_a\) be a unitary \(6\times6\) matrix whose columns form an
orthonormal basis of \(\mathbb C^6\). Mutual unbiasedness means

\[
|(B_a^*B_b)_{ij}|^2=\frac16\qquad(a\ne b).
\]

The unnormalized transition \(H_{ab}=\sqrt6 B_a^*B_b\) then has
unit-modulus entries and \(H_{ab}^*H_{ab}=6I\). All computational checks
must retain this normalization; testing an unnormalized product against
the normalized condition changes the problem. A complete family would
have seven bases and \(\binom72=21\) unordered transitions. A
fourth-basis obstruction and a seven-basis obstruction are distinct
statements.

\nolinkurl{MubQuartetTransportBridge.lean} now derives this
normalization for every actual quartet.
\nolinkurl{exists_quartet_iff_exists_rawQuartet} connects the existing
normalized matrix predicate to three raw Hadamard matrices with all
three MU conditions. \nolinkurl{MubMonomialTransport.lean} derives
preservation under a common left row-permutation/phase action and
independent right column actions. Thus standard Hadamard equivalence
preserves the entire completion problem.
\nolinkurl{no_quartet_of_concrete_representatives} still requires both
complete representative coverage and nonextension for every
representative. The separate pinned
\nolinkurl{MubFiniteCornerReduction.lean} now consumes the verified
external coverage theorem and derives an equivalence between actual
quartet existence and raw quartet completion of some retained
finite-corner output. Its nonexistence equivalence leaves the exclusion
of every retained output as the remaining mathematical statement, rather
than assuming coverage.

The Lean type \nolinkurl{MUB6.Gluing.Basis} and its MU relation are
opaque interfaces. Their names do not supply coordinates, unitary
matrices, or the above equations. A theorem over those interfaces needs
a separate concrete interpretation before it can be applied to the
Hilbert-space problem.

\subsubsection{1.1. Gauge coordinates for a fixed first Hadamard
matrix}\label{gauge-coordinates-for-a-fixed-first-hadamard-matrix}

Apply the same left unitary \(B_0^*\) to every basis. Then \(B_0=I\) and
\(H_a=\sqrt6 B_a\), for \(a=1,2,3\), have unit entries in modulus and
\(H_a^*H_a=6I\). A common left diagonal phase action, followed by
independent column phases, dephases \(H_1\). The common left action must
be applied to all bases; independently changing rows of just one basis
need not preserve MU.

For a fixed \(H_1\), right multiplication of \(H_2,H_3\) by diagonal
matrices with entries \(\overline{(H_a)_{0j}}\) fixes their first rows
to one. Entry moduli, unitarity and all transition moduli are preserved.
The acting column-phase group is \(\mathbb T^{12}\), of real dimension
\textbf{12}. The prescription fixes 24 ambient real coordinates: twelve
phases and twelve already-constrained radial coordinates. Thus
\(144-24=120\) real coordinates remain for the two additional raw
matrices. This is an ambient coordinate count, not a dimension theorem
for the solution set, a smoothness assertion, or a Jacobian rank claim.
There are still sixty circle equations before imposing orthogonality and
MU. If \(H_1\) also varies, its dephased \(5\times5\) block contributes
another fifty ambient real coordinates; the fixed-\(H_1\) count must not
be used globally.

The matrix-level preservation statements are
\nolinkurl{SzollosiColumnGauge.columnRephase_gram},
\nolinkurl{columnRephase_unitary}, \nolinkurl{columnRephase_rawPairMU}
and \nolinkurl{rawCandidates_firstRowDephase}. The family-specific
existence equivalence is
\nolinkurl{exists_gauged_solution_with_parameters_iff}; it keeps the
family parameters fixed and both additional bases variable.

\textbf{Scope.} This supports the stated gauge reduction. It does not
classify \(H_1\), establish a 120-dimensional manifold, or exclude a
quartet. The narrow claim is preservation of the completion problem
under column phases. The next required evidence is a complete
obstruction for the remaining equations over a declared domain,
including the parameters when they vary.

\subsubsection{1.2. Explicit transition-balance
identities}\label{explicit-transition-balance-identities}

Write \(A=H_2\), \(B=H_3\), \(W=A^*B\) and define polynomial residuals

\[
g_{Aj}=A_{0j}-1,\quad g_{Bj}=B_{0j}-1,\qquad
R_{Bm}=\sum_k B_{mk}\overline{B_{0k}}-6\delta_{m0}.
\]

For arbitrary complex matrices \(A,B\), without assuming any residual
vanishes, the following identity holds:

\[
\sum_k W_{jk}-6
=\sum_m\overline{A_{mj}}R_{Bm}
-\sum_{m,k}\overline{A_{mj}}B_{mk}\overline{g_{Bk}}
+6\overline{g_{Aj}}. \tag{1}
\]

Indeed, substituting \(\overline{B_{0k}}=1+\overline{g_{Bk}}\) cancels
the double sums on the right, leaving
\(\sum_{m,k}\overline{A_{mj}}B_{mk}-6\overline{A_{0j}}
+6(\overline{A_{0j}}-1)\). Interchanging \(A,B\) and conjugating gives
the column identity. On the first-row gauge and row-unitarity equations
both residuals vanish, so all six row sums and six column sums of \(W\)
are six. These are twelve complex, or twenty-four real, equations.

Equation (1), split into real and imaginary parts, is an explicit
polynomial ideal-membership identity for the row-unitarity and gauge
residuals. Its coefficients depend on matrix entries. For a square
matrix, column unitarity also implies row unitarity, as formalized by
\nolinkurl{SzollosiGaugeBalance.raw_right_gram}. The solution-level
conclusions are \nolinkurl{transition_row_sum},
\nolinkurl{transition_column_sum} and
\nolinkurl{solution_balance_equations}. No additional solution-set
restriction follows by adding these consequences, although they can
strengthen an interval contractor on boxes containing nonsolutions.

\textbf{Scope.} \nolinkurl{MubGaugeSyzygy.lean} now checks the row
identity and its conjugate column counterpart over every commutative
star ring. Its two ideal-membership corollaries apply to any ideal
containing the displayed generators, including polynomial ideals. They
support exact algebraic redundancy, without supplying a Jacobian rank or
nonextension result. All six displayed theorem endpoints were replayed
on 2026-10-11 with only standard Lean axioms. The next step is to bind
these generator identities to the executable contractor's coordinates.

\subsubsection{1.3. A quantitative invertible
minor}\label{a-quantitative-invertible-minor}

For any \(A\in\mathbb C^{3\times6}\) with \(AA^*=6I_3\), the finite
Cauchy--Binet identity gives

\[
\sum_{S\subset\{0,\ldots,5\},\ |S|=3}|\det A_{[:,S]}|^2
=\det(AA^*)=216.
\]

There are twenty summands. Hence one selected minor \(B=A_{[:,S]}\)
satisfies \(|\det B|^2\ge216/20=54/5\), in particular \(\det B\ne0\). If
all entries of \(A\) have unit modulus, each entry of
\(\operatorname{adj}B\) is a difference of two unit-modulus products.
Its squared modulus is at most four, and therefore

\[
|(B^{-1})_{ij}|^2\le\frac4{|\det B|^2}\le\frac{10}{27}.
\]

\nolinkurl{MubMinorChart.lean} checks the twenty-minor identity as an
exact polynomial identity and derives the two rational bounds. It also
applies them to the first three rows of every raw Hadamard matrix
satisfying \(HH^*=6I_6\). All ten displayed endpoints passed the
canonical Lean 4.33.1 replay with only standard axioms. The next
subsection supplies placement and complete reconstruction coverage.
Neither result establishes finite fibres or excludes any quartet. The
selected well-conditioned minor need not be the minor chosen by the
external finite-corner theorem.

\subsubsection{1.4. A complete bounded chart and division-free quartet
equations}\label{a-complete-bounded-chart-and-division-free-quartet-equations}

\nolinkurl{MubBoundedBlockChart.lean} extends each selected triple to a
complete column permutation, then dephases the entire matrix while
preserving determinant moduli and quartet extendability. On the block
index set \(\operatorname{Fin}(3)+\operatorname{Fin}(3)\) the resulting
matrix is

\[
H=\begin{pmatrix}E&B\\C&D\end{pmatrix},\qquad
q=|\det B|^2\ge54/5,\qquad D=-CE^*(B^{-1})^*.
\]

\nolinkurl{Admissible E B C} requires unit entries in all four blocks,
first row and column one in \(E\), first row one in \(B\), first column
one in \(C\), and the two diagonal row-Gram equations \(EE^*+BB^*=6I_3\)
and \(CC^*+DD^*=6I_3\). The forced formula makes the mixed Gram blocks
zero. The module checks soundness, coverage and actual quartet
transport. It does not use the external classification or a finite-fibre
premise.

\nolinkurl{MubBoundedQuartetEquations.lean} defines

\[
d=\overline{\det B},\qquad N=-CE^*(\operatorname{adj}B)^*,\qquad
S=\begin{pmatrix}dE&dB\\dC&N\end{pmatrix}.
\]

It shows in this model \(dD=N\) and \(dH=S\) under the positive
determinant bound. The forced-block constraints become \(|N_{ij}|^2=q\)
and \(qCC^*+NN^*=6qI_3\), without division. Retain the top Gram
equation, the unit entries of \(E,B,C\), their gauge constants and
\(q\ge54/5\). For two further raw matrices \(U,V\), retain unit entries
and \(U^*U=V^*V=6I_6\). Independent column phases fix their first rows
to one. The three MU conditions are exactly

\[
|(S^*U)_{ij}|^2=6q,\qquad |(S^*V)_{ij}|^2=6q,\qquad
|(U^*V)_{ij}|^2=6.
\]

The checked endpoint \nolinkurl{exists_quartet_iff_equations} identifies
actual quartet existence with a solution of these equations and the
determinant inequality. The matching nonexistence equivalence is checked
too. There are 76 free complex entries after the gauge constants are
substituted, or 152 ambient real coordinates. Unit entries bound each
real and imaginary part by one; this is not a solution-dimension or
independence claim. The scalar bound and all 34 displayed endpoints
across these two modules passed targeted canonical replay with only
standard axioms.

\nolinkurl{MubBoundedCoordinates.lean} checks both directions of the
152-real-coordinate codec and retains the full box \([-1,1]^{152}\),
including its boundary. \nolinkurl{MubBoundedPolynomials.lean}
represents complex arithmetic in the quadratic algebra over
\(\mathbb Q[x_1,\ldots,x_{152}]\) with generator squared equal to
\(-1\). Its real and imaginary components give 396 indexed rational
polynomial equalities. Together with \(q\ge54/5\), their real solutions
are equivalent to actual quartet existence
(\nolinkurl{exists_quartet_iff_polynomials}). The 13 coordinate and 26
polynomial endpoints passed scoped replay with only standard axioms.

\nolinkurl{MubBoundedExpressions.lean} identifies each executable
arithmetic expression with its polynomial, including determinant and
adjugate evaluation. The polynomials have total degree at most 8, and
\(q\) has degree at most 4.
\nolinkurl{MubBoundedCoefficientExpansion.lean} identifies finite
rational term-list expansion with those same polynomials. These degree
and expansion statements passed scoped kernel replay; they do not assert
sharp degree or independence.

\nolinkurl{MubBoundedExportIndices.lean} checks complete, unique
coordinate and equation numbering.
\nolinkurl{MubBoundedTermCollection.lean} checks polynomial preservation
under sorting, coefficient collection and exact-zero removal.
\nolinkurl{MubBoundedCodedSystem.exists_quartet_iff_coded} retains
actual quartet existence on \nolinkurl{Fin 152} with all 396 equations
and the determinant inequality. The 28 displayed endpoints passed scoped
replay with only standard axioms. The native export contains 3,760,338
collected monomials across the 396 constraints and \(q\); all 1,588
exact rational expression/list comparisons pass. The export receipt
binds the compressed artifact and its decompressed bytes.

\nolinkurl{MubBoundedMUFactor.lean} checks a further cancellation: 36 MU
residuals factor as \(q\) times polynomials of degree at most four.
Since \(q\ge54/5\), replacing those residuals preserves the complete
quartet existence equivalence. Their coefficient lists shrink from
371,520 to 3,096 monomials; with every other row retained, the total
becomes 3,391,914. The 36-row replacement table is bound to the original
archive by its receipt. Eighteen endpoints passed scoped replay. A
checked singular-point counterexample shows why the positive determinant
condition is required for this cancellation.

\nolinkurl{MubQuadraticLift.lean} supplies a separate full-matrix
formulation. Keeping the nine complex entries of \(D\) and three complex
transition matrices gives 386 ambient real coordinates.
\nolinkurl{MubQuadraticLiftPolynomials.lean} checks 648 rational
polynomial equalities: 108 unit-entry equations, 216 Gram components,
216 transition graph components and 108 transition norm equations. The
graph equations are \(W_0=H^*U\), \(W_1=H^*V\), \(W_2=U^*V\); each
transition entry has squared norm six. Actual quartet existence is
equivalent to a common zero in \([-1,1]^{170}\times[-3,3]^{216}\), with
coordinates grouped by their matrix-entry and transition slots.
\nolinkurl{MubQuadraticLiftDegree.lean} bounds the total degree of every
actual polynomial by two. The three modules' 28 displayed endpoints
passed scoped replay with only standard axioms.

This full-matrix model uses no determinant division. Its reverse
implication checks the three full Hadamard matrices and all pairwise MU
equations directly, so it does not need invertibility of \(B\). The
positive determinant condition remains necessary in the earlier
reconstruction and factor-cancellation models. The new coordinate and
equation counts are not independence or dimension claims.

\nolinkurl{MubQuadraticLiftIndices.lean} checks complete coordinate and
equation numbering and the bounds assigned to coordinate IDs. The
coefficient and coded-system modules identify the complete collected
lists with these same polynomials and preserve quartet existence on
\nolinkurl{Fin 386}. The native export contains 4,920 monomials, at most
12 per row, in 169,800 bytes. All 2,592 comparisons with direct rational
matrix arithmetic pass, and two exports are byte-identical. Forty-two
zero rows remain in the complete indexing.

\nolinkurl{MubQuadraticLiftTable.lean} contains literal rows generated
from that file. Its \nolinkurl{rows_match} checks every literal row
against the defining coefficient lists by kernel reduction;
\nolinkurl{bounds_match} checks the literal bounds.
\nolinkurl{exists_quartet_iff_table} preserves actual quartet existence
for common zeros of these literal rows in the coded box. JSON parsing
and file hashing remain external steps recorded in the coefficient
receipt. This does not identify an arbitrary external file with the
checked literals.

The native Rust loader now binds that exact file and checks its schema,
indexing, monomials and recorded probe values. One exact interval sweep
retains the first 170 coordinate bounds and narrows all 216 transition
bounds to \([-10035/4096,10035/4096]\). Every final equation range
contains zero; the result is \nolinkurl{CONTRACTED_UNRESOLVED_BOX}.
Three targeted Rust tests check this result, invalid-input rejection,
square enclosure, known-root retention and closed boundary contact.
These are native execution checks.

Independently, \nolinkurl{MubQuadraticLiftNarrowBox.lean} derives the
same bound from the transition norm equations. Its
\nolinkurl{exists_quartet_iff_narrow_table} retains actual quartet
existence on the smaller box. All seven displayed endpoints passed
scoped replay with only standard axioms. This checks the specific
mathematical narrowing without treating the runtime parser or general
contractor as kernel-reflected proofs.

\nolinkurl{MubQuadraticLiftCompatibility.lean} derives both transition
Gram matrices as \(36I\), the three cycle identities \(W_0W_2=6W_1\),
\(W_0^*W_1=6W_2\), \(W_1W_2^*=6W_0\), and row and column sums equal to
six. The polynomial companion adds 504 consequences: 216 column-Gram
components, 216 cycle components and 72 linear balance components. Their
degrees are checked, and adding them preserves actual quartet existence.
These new polynomials have not yet been exported as another literal
table.

The separation module checks, for every pair of columns,

\[
\sum_k |(W_0)_{ki}-(W_1)_{kj}|^2
=72-12\operatorname{Re}(W_2)_{ij}\ge\frac{43623}{1024}.
\]

The native inspector uses exact component interval bounds to exclude a
box when its squared-distance upper bound is strictly smaller. A closed
example with all coordinate intervals of positive width is excluded by
this test while the original 648 polynomial ranges still contain zero.
In each of two selected columns, real components lie in \([9/10,11/10]\)
and imaginary components lie within \(1/20\) of common alternating
centers \(9/4,-9/4\). The distance upper bound before contraction is
\(3/10\). \nolinkurl{close_columns_box_excluded} checks this family of
regions directly in Lean. Across the three new modules, 37 endpoints
passed scoped replay with only standard axioms. Five targeted native
tests passed, including the new regional regression and strict-threshold
contact. This supplies a regional exclusion; the full box remains
unresolved and no complete cover is claimed.

\textbf{Scope.} Infeasibility over the retained domain and a complete
exclusion cover remain open. The contraction and coefficient checks do
not establish emptiness or solver complexity; \(q\ge54/5\) remains part
of the older 152-coordinate system. The earlier large serialized archive
is separate from this newly reflected literal table.

The transition balances also permit a global affine reduction.
\nolinkurl{MubBalancedCompletion.lean} reconstructs a transition from
its leading \(5\times5\) block \(A\): the last column is
\(6-\sum_j A_{ij}\), the last row is \(6-\sum_i A_{ij}\), and the corner
is \(\sum_{i,j}A_{ij}-24\). Both reconstruction directions are checked
without division. \nolinkurl{MubBalancedQuadraticLift.lean} uses these
formulas for all three transitions, reducing 386 ambient real
coordinates to 320 while preserving actual quartet existence. The
retained box is \([-1,1]^{170}\times[-10035/4096,10035/4096]^{150}\).

The polynomial companion substitutes these affine formulas into all 648
original equations and checks that every resulting degree is at most
two. \nolinkurl{exists_quartet_iff_balanced_polynomials} is the
resulting global existence equivalence. The recovered coordinate bounds
hold on solutions; an arbitrary reduced-box point need not expand into
the old box. A checked counterexample uses the all-zero free assignment,
whose completed transition corner is \(-24\). All 27 displayed endpoints
across the three modules passed scoped replay with standard axioms.
Fewer coordinates do not establish a speedup or constraint independence.

The balanced coordinate codec and coefficient substitution now have a
complete exact export. The 18 displayed endpoints in
\nolinkurl{MubBalancedIndices.lean},
\nolinkurl{MubBalancedCoefficients.lean} and
\nolinkurl{MubBalancedCodedSystem.lean} passed scoped replay with
standard axioms, including \nolinkurl{exists_quartet_iff_coded}. The
generated \nolinkurl{MubBalancedTable.lean} contains literal equality
checks in blocks of 18 equations and a proposed
\nolinkurl{exists_quartet_iff_table} theorem. Its full build has not
completed. Its six displayed endpoints remain unverified; this
checkpoint does not claim a checked literal table. The exporter checks
2,592 rational values against direct completed-matrix arithmetic, and
two full exports have identical bytes. JSON parsing and byte custody
remain external to the kernel.

The balanced table contains 8,313 monomials, compared with 4,920 in the
386-coordinate table. Its largest row contains 676 terms rather than 12;
transition-norm equations account for most of the increase. The first
324 rows are unchanged. A corrupted corner row is rejected by the direct
matrix check, and the kernel separately checks that its true coefficient
list is nonempty. These are implementation and coefficient checks, not
exclusion certificates. The existing native inspector still uses 386
coordinates. Factored completion should be compared with expanded
substitution before choosing a reduced search representation; no runtime
or contraction-strength advantage has been established. Global
infeasibility remains open.

An independent Rust transcription in \nolinkurl{omega_mub} regenerates
all 648 rows of the balanced table from the same definitions and
compares them with the exported table term by term over the rationals;
every row agrees. This is a check of the export against its definitions,
not a replay of \nolinkurl{exists_quartet_iff_table}. Two negative
controls set the third basis equal to the first or to the second at the
product triple. In each case exactly the 36 transition-norm rows of the
repeated pair are nonzero, with values -6 and 30, and five core
coordinates leave the box. The Jacobian of the system at the centre of
the box has 1,045 linear terms in 269 rows and exact rank 184 of 320. It
is neither zero nor of full column rank, so a containment test whose
premise is a vanishing Jacobian at the centre does not apply to this
table, and a Moore--Penrose preconditioner is not defined there. The
receipt is \nolinkurl{artifacts/mub6_prompt05_campaign_receipt.json}.

\subsection{2. Exact arithmetic and finite
evidence}\label{exact-arithmetic-and-finite-evidence}

The computational fixtures use exact rational or cyclotomic arithmetic
over \(\mathbb Q(\zeta_{24})\). This coefficient field supports those
fixtures; it does not restrict arbitrary complex phases in the original
problem to 24th roots. Rational interval bounds can support continuum
statements only when they enclose the actual matrix-dependent functions
on a complete declared domain.

\nolinkurl{Conjecture2GeneralMechanization.lean} contains the finite
identities

\[
\alpha_2(m)=6m-9\in\{-9,-3,3,9\},\qquad m\in\{0,1,2,3\},
\]

\[
1+\operatorname{paritySign}(\alpha_2(m))=0,
\]

and paired integer cancellations. Its Tao phase-table calculation gives
the Eisenstein norm \(9\). The new \nolinkurl{balanced_mask_count} and
\nolinkurl{tao_all_balanced_masks} check all twenty
three-positive/three-negative sign patterns among the 64 six-bit masks.
For each column the signed exponent is reduced modulo three, and counts
\((n_0,n_1,n_2)\) give \(x=n_0-n_2\), \(y=n_1-n_2\), with exact norm
\(x^2-xy+y^2=9\). \nolinkurl{tao_distinct_positive_rows} also checks
every triple of distinct positive row indices directly, without a
bit-mask decoding premise. These four theorems were replayed using Lean
4.33.1 on 2026-10-11; \nolinkurl{#print axioms} reported no axioms for
them. \nolinkurl{MubTaoLaurent.lean} now supplies the complex
interpretation with \(\omega=-1/2+i\sqrt3/2\). It checks addition,
multiplication, conjugation and norm under the integer-pair evaluator,
identifies inverse entries with their conjugates, and derives

\[
\left|\sum_{j=0}^5\prod_{r=0}^5 S_{rj}^{\gamma_r}\right|^2=9
\]

for every balanced bit-mask sign vector \(\gamma\). The eight displayed
endpoints passed fresh replay using only standard Lean axioms; the
integer pair/exponent norm comparison uses none. This is a complex
matrix evaluation, not a necessary condition for arbitrary triples. The
original module's \nolinkurl{ModuliClass} table is a declared finite
table, rather than a classification of complex matrices. Its
analytic-certificate structure supplies row-count bounds and table
compatibility as assumptions. The later odd-parity polynomial lemmas
remain integer identities until a matrix monomial interpretation has
been constructed.

The SymPy 2-circulant calculations and Rust non-tensor surveys are
separate family evidence. They do not make every complex Hadamard matrix
a member of a surveyed family, or establish general Hadamard-cube
Conjecture 2.

\textbf{Scope of the Tao calculation.} It supports the exact nonzero
value for the recorded matrix and all balanced sign patterns. It does
not turn the conjectured Laurent vanishing condition for arbitrary MUB
triples into a theorem. The narrow claim is a finite exact invariant
evaluation. A Laurent-based exclusion requires that general
necessary-condition theorem, separately from the evaluation.

\subsubsection{2.1. Literature scope and classification
obligations}\label{literature-scope-and-classification-obligations}

The Fourier family is continuous, so it should not be grouped as a
sporadic matrix. \href{https://arxiv.org/abs/0902.0882v2}{Jaming et al.,
arXiv:0902.0882v2} {[}1{]} establish nonextension to a quartet for the
identity together with any member of their Fourier family. Applying that
result to an arbitrary transition requires common-unitary and monomial
equivalence transport; the new concrete transport modules supply
preservation, while matching the exact Fourier representatives to the
cited theorem remains a separate obligation.

\href{https://arxiv.org/abs/1402.4070v2}{Maxwell and Brierley,
arXiv:1402.4070v2} {[}2{]} treat the three-parameter Karlsson family and
establish its balanced Laurent vanishing identity. This does not give
quartet nonextension. The inherited description of the full Karlsson
family as a two-parameter affine family is not used here. Nor is an
unspecified matrix named \(H_6\) enough to identify a separate Haagerup
exclusion theorem: its concrete entries and source theorem must first be
matched.

\href{https://arxiv.org/abs/2608.18053v2}{Cárdenes Wuttig and Tindall,
arXiv:2608.18053v2} {[}3{]} present a finite-corner classification and
an accompanying Lean formalization. Their result concerns single
Hadamard matrices; it is not the five-family partition asserted in the
inherited campaign. The exact theorem and its equivalence conventions
were inspected against the concrete matrix predicates. The source was
pinned at \nolinkurl{57b03025ecbe259e474281f49e201ea7cb474e34}; its 136
project modules passed the source-level unfinished-proof/added-axiom
scan. After an initial memory failure and a successful single-worker
retry, the public theorem target completed under its pinned Lean
4.33.0-rc2 and Mathlib revision. All nine public theorem reports and the
supporting Karlsson coverage report use only standard axioms. All 215
external archive files matched the pinned ZIP.

The integration replay then recompiled the exact source bytes of 21
local dependency modules and the new bridge under the external pin. It
reconciles \nolinkurl{Fin 3 + Fin 3} with \nolinkurl{Fin 6}, row and
column Gram conventions, and actual monomial equivalence. All five
bridge endpoint reports use only standard axioms. The checked conclusion
is

\[
\begin{aligned}
&\exists\text{ an actual quartet in }\mathbb C^6\\
&\quad\Longleftrightarrow
\exists p,B,C:\ \operatorname{Retained}(p,B,C)
\land\operatorname{HasRawQuartet}(\operatorname{completion}(p,B,C)).
\end{aligned}
\]

Here \(p\) consists of four unit phases in a dephased corner \(E\); the
adjacent blocks obey \(EE^*+BB^*=6I\) and \(E^*E+C^*C=6I\), and the
forced final block is \(D=-CE^*(B^{-1})^*\). Retention includes the
external finite-corner and unit-entry conditions. Finite candidate
fibres over continuous seeds do not give a finite global list. Exclusion
over every retained reconstruction is still open. The old five-family
premise is quarantined.

\href{https://doi.org/10.1007/s10801-026-01506-x}{Matolcsi et al.,
Proposition 4.2} {[}4{]} provide a separate conditional route through
their Conjectures 2 and 3. Conjecture 2 concerns Laurent conditions on
actual MUB triples; Conjecture 3 concerns quartet nonextension of their
Szöllősi family. The paper explicitly withdraws an earlier treatment of
the latter as a literature theorem. Thus the local finite class table
cannot discharge either premise.

\textbf{Scope.} These source comparisons distinguish established family
results, replayed classification, concrete transport, and still-needed
exclusion arguments. They do not establish \(N(6)=3\). The narrow
conclusion is that the inherited five-branch case split cannot serve as
the current proof boundary. The next test is uniform exclusion on the
now-justified retained-output domain, or a direct proof of the two
conjectures in the alternative route. The independent quantitative minor
theorem does not identify its selected minor with the external
finite-fibre minor.

\subsection{3. Gluing statements and the Route C
frontier}\label{gluing-statements-and-the-route-c-frontier}

The combinatorial theorem \nolinkurl{g3_theorem} extracts a fourth basis
for any constituent triple of a supplied seven-basis family. The
assembly theorem combines supplied landing and obstruction data to
contradict that extension. This logical assembly does not construct the
landing data.

\nolinkurl{MUB6OrbitConfiguration} assumes

\[
6=n_D+n_T+n_S+n_H,\qquad n_T+n_S\le1,\qquad n_H\le0.
\]

These natural-number assumptions yield \(n_D\ge5\) and
\(n_D+\binom{n_D}{2}\ge15\). The structure contains no matrices or map
from these counts to actual transition edges. In particular, counting
pairs of labels does not establish that their relative transition
matrices lie in the Fourier family. The computation \(\min(2+1,3+1)=3\)
is also a numerical evaluation; it supplies no classification of
arbitrary complex MU bases.

The existing Route C conclusion has the exact premise

\begin{Shaded}
\begin{Highlighting}[]
\NormalTok{h\_landing : CompleteMUB6 → RouteCCompleteMUBFamily}
\end{Highlighting}
\end{Shaded}

The smaller interface \nolinkurl{FourierTripleLandingPremise} removes
orbit counts from the contradiction:

\[
\begin{aligned}
&\forall B:\operatorname{Fin}(7)\to\operatorname{Basis},\quad
\operatorname{PairwiseMU}(B)\Rightarrow\\
&\qquad\exists T:\operatorname{MUTriple},\quad
\operatorname{InConfig}(T,B)\land\operatorname{InFourierDitaChart}(T).
\end{aligned}
\]

\nolinkurl{no_complete_mub6_of_fourier_triple_landing} derives the
abstract obstruction from this supplied premise.
\nolinkurl{InFourierDitaChart} requires certified domain membership,
including analytic soundness. A concrete algebraic landing theorem and a
concrete interpretation of the chart remain open.

The proposed expression \(\sum_{a,b}|\operatorname{Tr}(B_a^*B_b)|^4\)
depends on basis-column phases. By contrast, the entrywise fourth-power
sum is fixed by MU: each distinct ordered pair contributes
\(36(1/6)^2=1\), and a self-pair contributes \(6\). For seven bases its
total is \(7\cdot6+7\cdot6=84\). This computation supplies no Fourier
membership map. A separate exact fixture checks that the ordinary raw
Fourier matrix has a \(2\times2\) cross ratio \(\zeta_6\) with
\(\zeta_6^4\ne1\). Row/column rephasing preserves cross ratios. Thus the
proposed all-fourth-root entry condition is not a definition of this
Fourier family. Likewise, a simplicial coboundary identity over integer
cochains contains no complex transition matrix data unless an additional
bridge has been established.

\href{https://arxiv.org/abs/0902.0882}{Jaming et al., arXiv:0902.0882}
{[}1{]} give a Fourier-family quartet obstruction. Applying that
literature result to an arbitrary seven-basis family still requires a
Fourier-family landing statement. It is not an unconditional
classification supplied by the citation.

\subsection{4. The 56-cell Krawczyk and reflection
frontier}\label{the-56-cell-krawczyk-and-reflection-frontier}

\nolinkurl{canonicalModuliPartitionCover56} stores 56 intervals with
endpoints \(i/56\) and \((i+1)/56\), a depth label \(64\), and fixed
positive rational margins. These data do not record a depth-64
contractor execution or the three cusp collars.

The kernel-checked equivalence
\nolinkurl{certified_point_in_cover56_iff} exposes the actual
hypothesis:

\[
\begin{aligned}
&\operatorname{CertifiedPointInCover56}(\mathrm{cover},\mathrm{pt})\\
&\quad\Longleftrightarrow
\operatorname{GeometricPointInCover56}(\mathrm{cover},\mathrm{pt})\\
&\qquad\land\operatorname{HasCurvatureCertificate}(\mathrm{pt.triple}).
\end{aligned}
\]

Thus \nolinkurl{conjecture3_szollosi_deltoid_closure} consumes a
curvature certificate; it does not derive that certificate from interval
containment.

The native library \nolinkurl{src/omega_native/mub6_smale_sentry.rs}
contains rational linear algebra and a Krawczyk contractor for supplied
square quadratic systems. The same-named binary under
\nolinkurl{src/bin/} now provides \nolinkurl{--inspect-system} for the
exported joint system. This mode performs exact interval inspection,
without a depth-64 contractor or a complete cell/cusp run. The terminal
receipt has aggregate labels and counts, without 56 per-cell rational
boxes, interval Jacobians, or Krawczyk images. Its status label is not
an independent replay of those missing objects.

A valid completion requires a concrete chart/system encoding, all cell
and collar domains, exact Jacobian bounds, complete coverage, and
checked exclusion on each leaf. Krawczyk strict inclusion supports root
existence under its hypotheses; exclusion instead requires disjointness
or another sound infeasibility condition. A \nolinkurl{by rfl}
evaluation of a Boolean establishes that computation. A separate
soundness theorem must connect it to the concrete MU residuals and
domain.

\nolinkurl{F_system} in the current abstract Lean module repeats a
basis-level MU predicate for the row and column indices. It is not yet
the 108 scalar coordinate polynomials needed by the proposed
native-to-kernel reflection. The cross-basis residuals also need the
candidate's unitarity/domain conditions: a root on an unrestricted
column array does not by itself represent a fourth basis.

The source family also needs careful domain identification.
\href{https://arxiv.org/pdf/0811.3930v2}{Szöllősi, arXiv:0811.3930v2}
{[}5{]} describes \(X_6(\alpha)\) as a two-real-parameter family over an
intersection of deltoid regions, containing the one-parameter
\(B_6(\theta)\) family. The stored scalar interval cells give no
coverage map for that full parameter domain. A native certificate must
specify the intended family and cover the associated domain.

\nolinkurl{SzollosiTwoCirculant.lean} now checks the continuous block
construction \(H=\begin{pmatrix}A&B\\B^*&-A^*\end{pmatrix}\), where the
circulant first rows are \((1,x,xy)\) and \((1,u,uv)\). Four unit phases
satisfy \(\phi(x,y)+\phi(u,v)=0\), with \(\phi(x,y)=x+y+\overline{xy}\).
The Gram identity is \(HH^*=6I\), so \(H/\sqrt6\) is unitary and all its
entry moduli squared are one sixth. The parameter domain is exactly six
scalar real polynomial equations in eight real coordinates. Its
solutions lie in \([-1,1]^8\), and
\nolinkurl{parameter_covers_intersection} connects those solutions to
every point of the intersection of the two phase-sum images. The cubic
root identities hold without a distinct-root premise. This is the source
block construction; equivalence transport to a selected dephased
representative and a classification of arbitrary Hadamard matrices are
not supplied.

\href{https://arxiv.org/pdf/2503.14752}{Matolcsi et al.,
arXiv:2503.14752, Conjecture 3} {[}4{]} concerns arbitrary quadruplets
containing a member of this family. A single chosen Zauner triplet would
not cover all possible third bases.
\nolinkurl{SzollosiQuartetSystem.lean} therefore leaves both additional
complex bases variable. \nolinkurl{quartet_iff_candidate_equations}
retains all 144 of their real coordinates and the actual pairwise MU
equations. Every coordinate of a matrix unbiased to the identity lies in
the rational box \([-1/2,1/2]\) by
\nolinkurl{mu_identity_coordinates_bounded}. The remaining empty-variety
obligation quantifies over every unitary third basis unbiased to both
fixed bases. No such variety is established empty.
\nolinkurl{SzollosiQuadraticSystem.lean} now gives a joint rational
polynomial specification: 440 real variables and 618 equations, each of
total degree at most two (\nolinkurl{equation_degree}). Raw matrices are
scaled by \(\sqrt6\) relative to normalized bases; three auxiliary
matrices represent their cross products, with entry norm-square target
six. The equation counts are:

\begin{longtable}[]{@{}lr@{}}
\toprule\noalign{}
Equations & Count \\
\midrule\noalign{}
\endhead
\bottomrule\noalign{}
\endlastfoot
Unit phases and complex compatibility & 6 \\
Family-matrix ties & 72 \\
Raw unitarity of the two additional bases & 144 \\
Unit entry norms of the two additional raw bases & 72 \\
Three cross-product matrix ties & 216 \\
Cross-product entry norms & 108 \\
\end{longtable}

\nolinkurl{exists_solution_with_parameters_iff} checks both directions
between actual quartets and solutions at every fixed admissible phase
parameter. \nolinkurl{exists_solution_iff_quartet} gives the full
existence equivalence. \nolinkurl{solution_in_full_box} bounds every
phase and raw basis coordinate in \([-1,1]\) and every auxiliary
cross-product coordinate in \([-3,3]\).
\nolinkurl{exists_bounded_solution_iff_quartet} checks that this
rational box loses no solutions. Both additional bases remain arbitrary,
including complex phases.

\nolinkurl{SzollosiExecutableSystem.lean} adds executable expression
trees for every checked equation. \nolinkurl{expression_polynomial} and
\nolinkurl{equation_eval} identify their real evaluation with the
rational specification for every assignment. \nolinkurl{variable_codes}
and \nolinkurl{equation_codes} check complete ordered IDs for all 440
coordinates and 618 equations. The exporter writes the full rational
containing box and three rational probes to
\nolinkurl{mub_szollosi_executable_system.json}; its raw SHA-256 is
\nolinkurl{1367d390da36e5ccf65abbc8fbf529ff07f9e154c277bb14c1cd8c36ba24c947}.
Independent exports were byte-identical.

The native inspector compiled 5,128 sparse terms of degree at most two,
matched all 1,854 probe values, and emitted all 618 full-box residual
intervals in \nolinkurl{mub_szollosi_native_replay_receipt.json}. Its
observed status is \nolinkurl{UNRESOLVED_FULL_BOX}. Finite probes do not
establish native symbolic parity for every real assignment.

\nolinkurl{SzollosiTermExpansion.lean} derives complete sparse
normalization semantics. The native inspector now compares all 618 rows
and 5,128 nonzero terms against the independently exported Lean
coefficient table, raw SHA-256
\nolinkurl{44656965bb72e735febfc69c1bcd599adbeea2e8512ebd6141025f8524e5d2df}.
\nolinkurl{SzollosiNativeReflection.lean} checks the actual emitted
coefficient table by kernel reduction (\nolinkurl{native_terms_match});
\nolinkurl{native_coefficients_polynomial} identifies every row with the
actual rational polynomial under injective coding.
\nolinkurl{SzollosiNativeInterval.lean} derives sparse interval
enclosure with square refinement and ordered accumulation. Kernel
reduction checks all 618 actual emitted intervals
(\nolinkurl{native_ranges_match}); \nolinkurl{native_interval_enclosed}
binds each to all actual real assignments in the full containing box.
The native receipt raw SHA-256 is
\nolinkurl{17b31a7eec56d0129b0b4b8a837f55d71b37ee0be0d0a91ad6fbb106949633dc}.
These identities concern the current emitted data, without generic
verification of every Rust compiler or JSON decoder input. Every emitted
interval still contains zero (\nolinkurl{native_full_box_unresolved}).
No variety emptiness or root existence follows. Coefficient and interval
literal mutations are rejected by their kernel checks; a quadratic
perturbation invisible to all three old probes is rejected by the
complete native coefficient comparison.

\nolinkurl{SzollosiIntervalExclusion.lean} derives rational AST interval
enclosure for every contained real assignment.
\nolinkurl{exclusion_check_sound} connects a successful residual
interval check to absence of actual polynomial zeros in that box. The
full containing-box check reduces to false inside the kernel
(\nolinkurl{full_box_check_unresolved}). A concrete first-phase box with
both coordinates in \([-1/2,1/2]\) is excluded by
\nolinkurl{phase_interior_excluded} (kernel reduction); this box
contains no admissible unit phase and supplies no exclusion of
admissible quartets. The complete containing box remains unresolved in
both calculations. A false exclusion check supplies no root-existence
conclusion.

\nolinkurl{RectangularQuadraticExclusion.lean} now supplies a rational
rectangular Krawczyk exclusion method.
\nolinkurl{list_derivative_polynomial} identifies the computed
derivative expressions with formal partial derivatives. The exact
quadratic midpoint secant identity yields \nolinkurl{root_in_image}:
every zero in the input box lies in its image for any rational
preconditioner. No inverse or square projection is required. Strict
separation in one coordinate gives \nolinkurl{exclusion_sound}; endpoint
contact is retained. A checked three-equation, two-coordinate example
separates a coordinate even though every natural residual interval
contains zero. A second checked example has strict image inclusion and
no common root, so strict inclusion alone supplies no root conclusion
for an overdetermined system.

\nolinkurl{SzollosiRectangularExclusion.lean} derives coefficient shape
from the checked native table and retains all 618 equations and 440
coordinates. \nolinkurl{derivative_enclosed} encloses formal derivatives
of the actual coded quartet polynomials.
\nolinkurl{no_solution_in_box_of_check} applies each successful leaf
check to actual polynomial solutions; \nolinkurl{no_quartet_of_check}
applies a successful containing- box check to actual quartets. Both
require an actual successful rational check and a center inside the box.
No successful quartet-domain exclusion or complete 56-cell/cusp run is
supplied.

The native rectangular command now emits the full containing-box
request, all 618 sparse Jacobian rows (9,000 entries), 440 images and
intersections. Its deterministic preconditioner has 28 nonzero rational
entries. The zero center lies in the containing box and need not satisfy
the phase equations. The output in
\nolinkurl{artifacts/mub_szollosi_rectangular_native_receipt.json} (raw
SHA-256
\nolinkurl{eccd3b1c075d0058cf729e2aa0af5e210860da93d2178ea09a9d5cfb3b820a76})
is \nolinkurl{CONTRACTED_UNRESOLVED_BOX}: 24 coordinates contract and
none separates. These contractions restore constant entries of the
constructed family matrix; none of the 144 raw coordinates of the two
additional bases is narrowed by this run, as checked by
\nolinkurl{native_additional_bases_unchanged}. The saved request
reproduces the receipt byte for byte.

\nolinkurl{SparseRectangularExclusion.lean} derives finite-support image
soundness. \nolinkurl{SzollosiRectangularReflection.lean} identifies
every emitted derivative interval and image with the actual polynomial
definitions, using bounded row checks evaluated inside the Lean kernel.
\nolinkurl{native_jacobian_entry_enclosed} bounds actual formal
derivatives. \nolinkurl{solution_in_native_retained} retains every
actual polynomial solution in every recorded intersection. The input
box, zero center, center residuals, contracted-coordinate list and
non-separation are also kernel checked. These checks apply to the
current emitted data; generic Rust/JSON verification is not supplied. No
root existence follows. \nolinkurl{SzollosiColumnGauge.lean} now derives
an independent column-phase reduction for both additional bases. It
leaves the constructed family matrix and all eight phase coordinates
fixed, and preserves the raw unitary, unit-entry and all three MU
Gram-norm equations. Conjugating each additional basis's first-row
entries sets that row to one. The same 618-equation system therefore has
a representative with 24 prescribed real coordinates whenever a quartet
exists, also at every fixed admissible family parameter. These
coordinates are wire IDs 80-91 and 152-163, with alternating real value
one and imaginary value zero; their complete ordered list is checked by
\nolinkurl{gaugePointCodes_checked} using \nolinkurl{by rfl}.
\nolinkurl{gauge_box_firstRow} derives the first-row equalities from
point-box membership;
\nolinkurl{exists_gauge_box_solution_with_parameters_iff} gives the
exact existence reduction with no separate gauge premise. The bounds are
exported in \nolinkurl{artifacts/mub_szollosi_column_gauge_box.json}.
This normalizes columns of the two variable bases; it does not classify
or fully dephase the family matrix. The earlier full-box contractor
receipt is unchanged. No quartet-domain exclusion or complete cell/cusp
coverage follows from this reduction.

\nolinkurl{SzollosiGaugeBalance.lean} derives row and column sums of six
for the raw transition between the two additional bases after the column
gauge. \nolinkurl{SzollosiGaugedExecutable.lean} retains all 618
original equations and appends 24 gauge equations and 24 derived balance
equations. The resulting 666 quadratics use the same 440 coordinates and
preserve quartet existence at every fixed admissible family parameter.
The balance equations are conclusions of the original equations and
gauge, rather than new assumptions on quartets. The Lean export and
native replay agree on all 666 coefficient rows, 5,320 sparse terms and
1,998 rational probe values; the first 618 rows are unchanged.

\nolinkurl{SzollosiGaugedCover.lean} checks a six-chart solution cover:
at least one real entry in the first transition row is at least one,
since that row sums to six. The six charts overlap.
\nolinkurl{SzollosiGaugedReflection.lean} connects the actual native
chart bounds to every augmented-system solution, also at each admissible
family parameter, using the earlier checked retained-box theorem. Both
additional bases remain variable. Duplicate native runs on all six
charts return \nolinkurl{UNRESOLVED_BOX}; none of these charts has been
excluded.

A separate candidate box restricts all six real entries of that
transition row to {[}-1/2,1/2{]}. Its balance equation forces the first
real entry into {[}7/2,17/2{]}, separated from its own input interval.
The native receipt
\nolinkurl{artifacts/mub_szollosi_balance_exclusion_receipt.json}
records this image at coordinate 368. Kernel checks bind the 48 added
coefficient rows and all 440 input bounds; the selected image has
independent real semantics from the balance identity. A contained point
has admissible phases (1,-1,1,i), so this box's exclusion is not an
off-phase-domain example. That point is not a quartet. The result
excludes this candidate box, without excluding all candidates at any
family parameter. Other newly emitted images and the generic Rust/JSON
runtime are not verified by this selected-image result. The new root
replay fixes 48 coordinates (24 constructed-family constants and the 24
column gauges) with 74 nonzero preconditioner entries. It remains
unresolved. Six candidate charts do not replace the required 56 family
cells and three cusp collars, and \nolinkurl{HasCurvatureCertificate}
remains an explicit premise.

Useful preconditioners, further subdivision, further equivalence
transports, complete parameter/candidate coverage, cusp analysis and
curvature remain open. The native runtime receipt retains its
\nolinkurl{kernel_reflection_complete: false} field; the separate
checkpoint binds the subsequent Lean checks without altering the native
output.

The Lean \nolinkurl{SmalePointCertificate} and
\nolinkurl{KrawczykContractor} structures already contain a root, its
equation, and uniqueness evidence. The current existence lemmas project
these fields; \nolinkurl{krawczyk_inclusion_criterion} does not derive
the root from \nolinkurl{K_subset_interior}. They therefore supply no
independent analytic soundness bridge for a native contractor receipt.

The separate deltoid-cover implementation restricts candidates to
\nolinkurl{canonical_candidate_basis_box()}, whose first real component
lies in \([9/10,1]\). Consequently \(|U_{0,0}|^2 \ge 81/100\), and its
base margin is \((81/100-1/6)^2=37249/90000\). Every basis unbiased to
the standard basis has \(|U_{0,0}|^2=1/6\), outside that box. This
margin therefore excludes a restricted candidate domain that contains no
such unbiased basis. Subdividing the family parameter domain, even
exhaustively, does not enlarge the candidate domain. The restriction is
explicit in
\nolinkurl{mub6_szollosi_deltoid_cover.rs::canonical_candidate_basis_box}
and \nolinkurl{SzollosiDeltoidCoverReport.candidate_domain_restriction}.

\subsection{5. The flat-extension and operator
frontier}\label{the-flat-extension-and-operator-frontier}

For a genuine normalized moment functional \(L\), a degree-two moment
matrix has entries \(M_2(L)_{u,v}=L(u^*v)\) for a declared word basis. A
flat-extension bridge must bind \(M_1\) to the corresponding principal
submatrix, require positivity and relation compatibility, and establish
actual equality of ranks. The legacy rank labels in
\nolinkurl{RankLoopClosure} and \nolinkurl{CurtoFialkowFlatExtension}
are not computed ranks of these matrices. The added
\nolinkurl{rationalMomentMatrix} evaluates explicit word products in
\(\mathbb Q\). \nolinkurl{MomentPrincipalEmbedding} binds the lower word
basis to a principal submatrix, and
\nolinkurl{rationalMomentMatrix_rank_mono} establishes rank
monotonicity. \nolinkurl{FlatMomentRankCondition} now states equality of
the actual degree-two and degree-one matrix ranks, together with this
inclusion and complete, nonduplicate indexing of all words at both
degrees. Empty index lists cannot satisfy it. No MUB-specific flatness
witness, relation compatibility, or extraction follows from that
definition.

\nolinkurl{gns_extraction_theorem} takes a structure already carrying
the operators and state. It repackages that data into the present
realization interface. The legacy theorem supplies no operator
extraction from moment values. The rank-loop exclusion theorem retains
\nolinkurl{h_gns_exhaust}; the emptiness variants retain their soundness
or trace-sign arguments.

There is a preceding domain defect.
\nolinkurl{PhysicalProjectorRealization} does not require basis
completeness, rank-one ranges, Hilbert dimension six, inner-product
linearity, or cross-basis unbiasedness.
\nolinkurl{zeroActionInterfaceModel} supplies a Boolean carrier, zero
actions, and a normalized designated state satisfying every current
field. The theorem \nolinkurl{not_universal_interface_exclusion} then
gives

\begin{Shaded}
\begin{Highlighting}[]
\NormalTok{¬ UniversalPhysicalProjectorExclusion numBases dim}
\end{Highlighting}
\end{Shaded}

for arbitrary size labels, with no axiom dependencies. This is an
interface countermodel, not a concrete quantum or MUB realization. It
prevents treating the present exclusion predicate as the requested
physical theorem.

\nolinkurl{ConcreteMubProjectors.lean} provides a separate forward
construction from \nolinkurl{FiniteComplexMubFamily ι B ι}. Its actual
complex matrices are \(P_{b,l}=v_{b,l}v_{b,l}^*\). The checked lemmas
establish self-adjointness, same-basis multiplication,
\(\sum_l P_{b,l}=I\), trace one, rank one, and
\(\operatorname{Tr}(P_{b,l}P_{c,k})=1/|\iota|\) for distinct bases. This
constructs the projectors from a supplied complex family; it neither
supplies such a family in dimension six nor extracts it from positive
truncated moments. The legacy interface countermodel therefore remains
relevant.

\nolinkurl{ConcreteMubMoments.lean} now constructs complex word moments
directly from that supplied family. Choose a state \(\psi=v_{b_0,l_0}\)
and a distinct basis \(b_1\ne b_0\). For a word \(w\), let \(P_w\) be
the ordered projector product and
\(L_\psi(w)=\langle\psi,P_w\psi\rangle\). The complete, nonduplicate
word lists of lengths at most one and two index actual Hermitian
matrices. The theorem \nolinkurl{complexMomentMatrix_eq_gram} identifies
each matrix with the Gram matrix of its word-state columns;
\nolinkurl{complexMomentMatrix_posSemidef} establishes positivity.
\nolinkurl{complexMomentMatrix_principal_degree_one} identifies the
lower principal submatrix. MU makes every overlap with basis \(b_1\)
nonzero, so its singleton projector columns span the coordinate space.
Consequently \nolinkurl{concrete_mub_moment_flatness} gives
\(\operatorname{rank}(M_2)=\operatorname{rank}(M_1)=|\iota|\).
Normalization at the empty word is also checked.

This is a forward theorem for moments induced by a supplied concrete
family, with the distinct-basis and state labels explicit. It supplies
neither a family in dimension six nor a reverse extraction from
arbitrary positive flat data. The rational moment API remains separate;
an adequate MUB interpretation of abstract positive flat data remains
required.

\nolinkurl{ConcreteMubPolynomials.lean} now supplies actual coordinate
polynomials for a fixed triple of complex matrices. A candidate
\(U\in\mathbb C^{6\times6}\) is encoded by 72 real coordinates, with a
checked inverse decoding. The 108 MU residuals are
\(r_{a,i,j}=|\langle T_{a,i},U_j\rangle|^2-1/6\). The 72 real/imaginary
residuals of \(U^*U-I\) impose candidate orthonormality. These are 180
indexed constraints, without an independence claim.
\nolinkurl{polynomialIdeal} is their generated ideal in
\nolinkurl{MvPolynomial Coord} over \(\mathbb R\);
\nolinkurl{polynomialIdeal_le_eval_kernel_iff} identifies its
evaluation-kernel condition with their simultaneous vanishing.
\nolinkurl{muPolynomial_eval_projector_trace} identifies each MU
residual with the corresponding complex projector cross-trace minus one
sixth.

\nolinkurl{polynomialVariety_eq_empty_iff} connects this actual real
zero set to absence of a unitary candidate with the stated MU overlaps.
A supplied complete complex MUB family gives a point in the zero set of
its first three bases, using its fourth basis as candidate, through
\nolinkurl{family_triple_variety_nonempty}. This is a checked semantic
bridge. It does not establish zero-set emptiness. The separate
commutative construction below now extracts a candidate from normalized
positive flat coordinate moments satisfying the finite residual
equations. The fixed triple's coefficients are real and imaginary parts
of complex entries; rational enclosures over the continuous family cells
and cusp collars still require separate construction.

\nolinkurl{FlatMomentProjectorExtraction.lean} now gives a finite
reverse construction. Its inputs are normalized complex word values,
positivity of the actual complete degree-two moment matrix, equality
with the actual degree-one matrix rank, and contextual idempotence
equations \(L(u^*ggv)=L(u^*gv)\) for each generator and lower-degree
words. It constructs a Gram factor, one common complex linear space, a
normalized state, and linear idempotent actions symmetric for the
explicit finite Hermitian pairing. \nolinkurl{extract_word_projectors}
identifies every degree-two moment entry with
\(L(u^*v)=\langle P_u\psi,P_v\psi\rangle\), for words of length at most
two. The space dimension is \(\operatorname{rank}(M_1)\). Operators and
a state are outputs, rather than input structure fields.

\nolinkurl{concrete_moment_extraction_inputs} checks that the previously
constructed moments from a supplied MUB family satisfy these extraction
premises. \nolinkurl{FlatMomentBasisRelations.lean} adds the exact
contextual equations \(L(u^*g_{b,l}g_{b,k}v)=\delta_{lk}L(u^*g_{b,l}v)\)
and \(\sum_l L(u^*g_{b,l}v)=L(u^*v)\) for every basis and all words of
length at most one on either side. These equations are input relations,
rather than consequences of positivity alone.

From these relations, \nolinkurl{extract_word_basis} constructs actions
satisfying \(P_{b,l}P_{b,k}=\delta_{lk}P_{b,l}\) and \(\sum_lP_{b,l}=I\)
on the common extracted space, while retaining symmetry, state
normalization, and every degree-two moment entry.
\nolinkurl{concrete_moment_basis_inputs} checks the contextual relations
for moments of a supplied MUB family. The normalized word assignment
that is zero on every nonempty word fails completeness, as checked by
\nolinkurl{zero_action_moments_reject_completeness}.

For arbitrary input word data, dimension six, nonzero rank-one ranges
for every outcome, and cross-basis MU conclusions still need their
corresponding relations and derivations. The finite representation does
not assert a full star-algebra extension or transport arbitrary word
data to the coordinate polynomial zero set. These remain limits of that
noncommutative word route. The real coordinate route below supplies a
separate extraction and its empty-variety consequence without that
additional word-to-coordinate transport.

\nolinkurl{CommutativeFlatMoments.lean} constructs real coordinate
multiplication operators directly. Let the lower index list contain the
constant and every coordinate, and the upper list also contain every
ordered coordinate product. The repeated quadratic monomials are
retained. For a real linear functional \(\Lambda\), the actual entries
are \(M_2(\Lambda)_{i,j}=\Lambda(m_i m_j)\). Normalization
\(\Lambda(1)=1\), positive semidefiniteness and the actual equality
\(\operatorname{rank}M_2=\operatorname{rank}M_1\) are explicit premises.

A Gram factor and rank equality identify the full column space with its
lower column space. The coordinate shifts descend to symmetric
endomorphisms there; polynomial commutativity yields commuting
operators. Their joint real eigenspaces span this finite space. A joint
eigenvector with nonzero overlap with the normalized state yields one
common real point. Every quadratic polynomial \(p\) satisfying
\(\Lambda(p^2)=0\) vanishes at that point, through
\nolinkurl{real_point_of_positive_flat}. The operators and point are
constructed, rather than supplied as extraction fields.

\nolinkurl{ConcreteMubFlatExtension.lean} checks that all 180
candidate-basis residuals \(r_e\) lie in this quadratic span.
\nolinkurl{FiniteRelationCompatibility} states \(\Lambda(m_i r_e)=0\)
for every upper-list monomial and every residual. These finite equations
imply \(\Lambda(r_e^2)=0\), so the same extracted point annihilates all
residuals. \nolinkurl{fourth_basis_of_positive_flat} produces an actual
complex unitary candidate with all 108 MU overlaps for the fixed triple.

Evaluation at any actual solution supplies a normalized positive moment
matrix with both ranks equal to one and satisfies every finite relation
equation. \nolinkurl{fourth_basis_iff_positive_flat_moments} gives the
resulting two-way equivalence between candidate existence and existence
of normalized positive flat compatible coordinate moments.
\nolinkurl{moment_nonpositive_of_empty_variety} consequently excludes
positive semidefiniteness when the fixed real variety is empty,
retaining normalization, actual flatness and finite relation
compatibility as hypotheses. It does not exclude arbitrary nonflat
truncations or establish variety emptiness.

The concrete flat-moment, normalized-quartet, and joint polynomial
equivalences were freshly replayed as dependencies of the 2026-10-11
transport bridge, with only \nolinkurl{propext},
\nolinkurl{Classical.choice}, and \nolinkurl{Quot.sound} in their
reported axioms. This scoped replay covers 23 local source modules and
cached dependencies; it is not a fresh build of the complete MUB6 import
surface. This real coordinate model is separate from the rational and
noncommutative word APIs; no implicit identification of their data is
used.

\href{https://arxiv.org/abs/0812.2563}{Laurent and Mourrain,
arXiv:0812.2563} {[}6{]} and
\href{https://arxiv.org/abs/1406.4975}{Mourrain and Schmüdgen,
arXiv:1406.4975} {[}7{]} provide commutative and star-algebra
flat-extension settings, respectively. The checked finite commutative
construction above derives its own operator and point extraction from
explicit hypotheses. A full noncommutative extension is not inferred
from those references. Variety emptiness does not, by itself, exclude an
arbitrary fixed truncation; normalization, actual flatness and relation
compatibility remain part of the stated consequence.

\subsection{6. Receipt custody and
validation}\label{receipt-custody-and-validation}

The R6 receipt is present and has raw SHA-256
\nolinkurl{cf8442d8fe1397226d669fcea7bc038ef8d984a9bae741772b9d742fbe62eb59}.
It records \nolinkurl{NOT_A_CERTIFICATE}.

The Matolcsi receipt has raw SHA-256
\nolinkurl{598d73b4f7b6807d62fc8d5c7698fd7c00ca73634c24641e0f2a848c873fc6fa};
its internal digest is a different commitment. The historical
\nolinkurl{zenodo/release_manifest.sha256} contains 496 entries: 493
archive file members plus three release-metadata files. The archive
extraction regression checks all 493 file members. Digest agreement
establishes the corresponding byte or normalization commitment, rather
than theorem validity.

The research source snapshot is pinned at
commit
2d03b504106356593d76f5a6bafc259019d52247. The
Lean
source package and
Rust
workspace are external research materials. The preprint ZIP does not
bundle those trees, their dependencies or a new axiom census. Its
integrity check verifies only the delivered manuscript package. The
historical research archive and receipt counts above describe a separate
artifact, not the preprint ZIP.

\nolinkurl{Mub6} is the canonical package root; the external-pinned
integration target is separate. \nolinkurl{Mub6} currently imports
\nolinkurl{MubBalancedTable}, whose full build has not completed
(Section 1.4), so a successful \nolinkurl{lake build Mub6} has not been
recorded for this checkpoint. The full build and all-declaration axiom
audit must complete before a package-wide formal validation claim is
made. The live execution evidence and outstanding gates are recorded in
the executive checkpoint report; historical test totals are not
substituted for current runs. The 2026-10-11 finite-corner receipt
separately binds the canonical index/minor targets, the external public
theorem, the 22-source integration replay, and their printed axiom
reports. It does not assert a whole-package replay. The subsequent
bounded-quartet receipt binds the two canonical block-chart and
division-free equation targets, their 34 endpoint reports, and the
rejected determinant-conjugation mutation.

The bounded export receipt binds 31 local source modules, 28 displayed
index, collection and coded-system endpoints, two rejected mutations,
and the complete native coefficient artifact. Its source-level
equivalence is kernel-checked; serialized-byte reflection and global
infeasibility remain open. These scoped receipts do not report execution
of the complete validation command block above.

\subsection{7. Publication checkpoint and completion
requirements}\label{publication-checkpoint-and-completion-requirements}

The package is not ready for an arXiv submission asserting mathematical
closure, or for a new Zenodo release with that claim. A scoped technical
checkpoint must still complete rendering, metadata reconciliation,
semantic review, and archive re-certification before deposit. The
existing archive is historical and does not contain the current edits by
implication.

The full Option B objective is retained: establish a concrete Route C
landing bridge or state its defensible remaining premises; execute and
soundly reflect the complete 56-cell/cusp exclusion calculation;
construct the operator moment and flat-extension bridge in an adequate
mathematical domain. A passing suite or a standard-axiom implication
cannot replace any of those requirements.

\subsection{8. Mechanism inventory and evidence
boundaries}\label{mechanism-inventory-and-evidence-boundaries}

\begin{longtable}[]{@{}
  >{\raggedright\arraybackslash}p{(\linewidth - 4\tabcolsep) * \real{0.3333}}
  >{\raggedright\arraybackslash}p{(\linewidth - 4\tabcolsep) * \real{0.3333}}
  >{\raggedright\arraybackslash}p{(\linewidth - 4\tabcolsep) * \real{0.3333}}@{}}
\toprule\noalign{}
\begin{minipage}[b]{\linewidth}\raggedright
Mechanism
\end{minipage} & \begin{minipage}[b]{\linewidth}\raggedright
Mathematical object or conclusion
\end{minipage} & \begin{minipage}[b]{\linewidth}\raggedright
Evidence boundary
\end{minipage} \\
\midrule\noalign{}
\endhead
\bottomrule\noalign{}
\endlastfoot
Concrete normalization and transport & Raw and normalized quartet
equivalence & Reported external Lean endpoints; a current full replay
remains required \\
Quantitative minor and block chart & Squared determinant at least 54/5;
inverse-entry squared moduli at most 10/27 & Section 1.3 and named
external Lean endpoints; no nonextension conclusion \\
Balanced polynomial reduction & 648 equations in 320 bounded real
coordinates & Reported degree and equivalence endpoints; infeasibility
open \\
Fourier-family route & Conditional contradiction from a Fourier-triple
landing premise & The landing and concrete chart-soundness premises
remain open \\
Interval exclusion and six charts & Sound exclusion criteria and
recorded candidate boxes & No complete admissible-family exclusion or
56-cell/cusp replay \\
Positive flat moment extraction & Fourth-basis construction from
explicit compatible moment hypotheses & Empty varieties and flatness for
proposed data remain open \\
Finite Tao calculation & Norm 9 on all twenty balanced sign patterns &
Finite matrix evaluation; no general Laurent necessary-condition
theorem \\
Authorization & None & Integrity checks and reported computations confer
no mathematical or publication approval \\
\end{longtable}

The retained research ledgers are historical working indexes in the
source snapshot. They are not part of the standalone preprint or its
verification contract. The source theorem names in Sections 1--5 specify
the individual reported results; their stated hypotheses remain part of
each statement.

\subsection{References}\label{references}

\begin{enumerate}
\def\labelenumi{\arabic{enumi}.}
\tightlist
\item
  P. Jaming, M. Matolcsi, P. Móra, F. Szöllősi and M. Weiner, ``A
  generalized Pauli problem and an infinite family of MUB-triplets in
  dimension 6'',
  \href{https://arxiv.org/abs/0902.0882v2}{arXiv:0902.0882}.
\item
  A. Maxwell and S. Brierley, ``On properties of Karlsson Hadamards and
  sets of Mutually Unbiased Bases in dimension six'',
  \href{https://arxiv.org/abs/1402.4070v2}{arXiv:1402.4070}.
\item
  M. Cárdenes Wuttig and J. Tindall, ``A Complete Classification of
  Complex Hadamard Matrices of Order Six'',
  \href{https://arxiv.org/abs/2608.18053v2}{arXiv:2608.18053v2}.
\item
  M. Matolcsi, Á. K. Matszangosz, D. Varga and M. Weiner, ``Triplets of
  mutually unbiased bases'', Journal of Algebraic Combinatorics (2026),
  \href{https://doi.org/10.1007/s10801-026-01506-x}{doi:10.1007/s10801-026-01506-x};
  \href{https://arxiv.org/abs/2503.14752}{arXiv:2503.14752}. Cited in
  Sections 2.1 and 4 under both identifiers.
\item
  F. Szöllősi, ``A two-parameter family of complex Hadamard matrices of
  order 6 induced by hypocycloids'',
  \href{https://arxiv.org/abs/0811.3930v2}{arXiv:0811.3930v2}.
\item
  M. Laurent and B. Mourrain, ``A Sparse Flat Extension Theorem for
  Moment Matrices'',
  \href{https://arxiv.org/abs/0812.2563}{arXiv:0812.2563}.
\item
  B. Mourrain and K. Schmüdgen, ``Flat extensions in *-algebras'',
  \href{https://arxiv.org/abs/1406.4975}{arXiv:1406.4975}.
\end{enumerate}

\end{document}
